\subsection{Retaining a quarter of the second Frobenius block}
\label{sec:quarter-frobenius-block}

\begin{proposition}\label{prop:quarter-frobenius-block}
For every prime $p\ge257$, there is a dimension-three Reed--Solomon code
over $E=\mathbb F_{p^4}$ of length
\[
 n=\frac{5p^2-1}{2}
\]
and two received words $r,s$ with individual and common agreements all
exactly $A=2p$ such that, at threshold $T=\lceil\sqrt{5p^2-1}\rceil-1$, the
normalized affine line $(1-z)r+zs$ has exactly
\[
 B=(p+1)(p^2-1)
\]
parameters with singleton lists, one further parameter with list size
$p+1$, and empty lists at all other parameters. Moreover,
\[
 n a_1(3/n)<A<T<\sqrt{2n}.
\]
Thus $B=\Theta(n^{3/2})$, the relative source and common-agreement loss
$(T-A)/(T-3)$ tends to $1-2/\sqrt5$, and the characteristic is
$p=\Theta(\sqrt n)>2$. The threshold is the largest integer strictly below the exact Johnson
agreement bound. The evaluation domain is a selected subset of
$\mathbb F_{p^4}$; this is not a prime-alphabet construction.
\end{proposition}

The code encodes three field elements, the coefficients of a quadratic,
into $n$ evaluations. Its rate is $3/n=\Theta(p^{-2})$, and the
fractional agreement loss $(T-A)/n$ is $\Theta(p^{-1})$.
The constant ratio above is relative to the vanishing capacity margin
$(T-3)/n$, not a constant loss in relative distance.

\begin{proof}
Use $B_0=\mathbb F_{p^2}$, $\eta=s_1^2\notin B_0$, the two blocks
$D_0,D_1=s_1D_0$, and the received words $f,g$ from the proof of
Proposition~\ref{prop:scaled-frobenius-padding}, before its neutral
padding. Thus $|D_0|=|D_1|=2(p^2-1)$, and the quadratics canonical on
both blocks at parameter $\lambda$ are precisely
\[
 Q_{a,b}=aX^2+b,\qquad
 a^{p+1}=1,\quad b,v\in I_a,\quad\lambda=b-\eta v.
\]
Their first-block supports have size $2p$ or $2p-2$, according as
$b\ne0$ or $b=0$. Their second-block supports depend only on $(a,v)$
and have size $2p$ or $2p-2$, according as $v\ne0$ or $v=0$.
There are $p(p+1)$ such second-block supports.

There is a subset $D'_1\subset D_1$ of size
$m=(p^2+3)/2$ meeting every one of these supports in at least
$\lceil p/4\rceil$ points. Indeed, choose $D'_1$ uniformly among the
$m$-subsets. For a support of size $k\in\{2p-2,2p\}$, its retained
size $Z$ is hypergeometric with mean greater than $k/4$. Since
$\lceil p/4\rceil-1\le(p-1)/4\le k/8$, the usual Hoeffding bound
for sampling without replacement gives
\[
 \Pr[Z<\lceil p/4\rceil]\le \exp(-k/32)
 \le\exp(-(p-1)/16).
\]
Consequently the union bound is less than one for every $p\ge257$.
For an explicit onset check, $257\cdot258<16^6/6!+16^7/7!<e^{16}$,
and $\log(x(x+1))-(x-1)/16$ decreases for $x\ge257$.
Fix such a subset and use $D=D_0\cup D'_1$, with no neutral padding.

The quadratic classification in the cited proof also gives two bounds
that hold after arbitrary puncturing of the second block: a quadratic
canonical on only one block has at most $2p$ total matches, and a
quadratic canonical on neither block has at most $p+4\le2p$ matches.
Thus every witness above $A$ must be canonical on both blocks. Every
such witness still has at least
\[
 2p-2+\lceil p/4\rceil>\sqrt{5p^2-1}>T
\]
matches: indeed $2p-2+\lceil p/4\rceil\ge9p/4-2>0$, and
\[
 (9p/4-2)^2-(5p^2-1)=\frac{p^2-144p+80}{16}>0
\]
for $p\ge257$. The planes $I_a+\eta I_a$ intersect pairwise only at zero.
Each of their $B$ nonzero elements therefore has exactly one witness
at threshold $T$, while zero has exactly the $p+1$ witnesses $aX^2$.
No other parameter has a qualifying witness.

Choose distinct $\alpha,\beta$ outside the union of these planes, and
put $r=f+\alpha g$, $s=f+\beta g$. Such choices exist because the
union has $1+B<p^4-1$ elements. The classification bounds each source
agreement by $2p$, and any first-block canonical quadratic with
$b\ne0$ attains $2p$ matches simultaneously for the two sources on
$D_0$. Hence both source agreements and common agreement equal $A$.
The bijection $\lambda=\alpha+z(\beta-\alpha)$ preserves the complete
list profile and has no omitted or infinite parameter.

Finally, using the same first-order upper bound as above,
\[
 n a_1(3/n)\le\sqrt{3n/2}+(3n/8)^{1/4}
 <\frac{31}{16}p+\sqrt p<2p
\]
for $p>256$. Also $(2p+1)^2\le5p^2-2$ for $p\ge257$, so
$T=\lfloor\sqrt{5p^2-2}\rfloor\ge2p+1>A$, while
$T^2<5p^2-1=2n$ by its definition. This proves the stated comparisons.
\end{proof}

\begin{corollary}[Growing dimension, with weaker placement]
\label{cor:quarter-frobenius-growing-dimension}
The construction in Proposition~\ref{prop:quarter-frobenius-block} can
instead have message dimension $k=\lfloor p/20\rfloor+1$, at the same
length $n$ and threshold $T$, with exactly the same complete threshold-list
profile. Both source agreements and their common agreement then lie in
$[2p,2p+2\lfloor p/20\rfloor]$. The guaranteed loss, divided by $T-k$,
has lower limit at least
\[
 \frac{\sqrt5-21/10}{\sqrt5-1/20}>0.0622.
\]
The rate is now $\Theta(p^{-1})$, but the threshold is below the enlarged
code's first-order curve. This is not an improvement in the
above-first-order regime of the proposition.
\end{corollary}

\begin{proof}
Choose the scale $s_1$ in the proposition to be primitive in $E^*$.
Write $B_0=\mathbb F_{p^2}$ and choose $v\notin B_0$ with $v^2\in B_0$.
The two physical $B_0$-lines in each block are
$B_0,vB_0$ and $s_1B_0,s_1vB_0$, respectively, with zero removed.
Before puncturing, the received words are
\[
 f(X)=X^{2p},\quad g(X)=0\quad\text{on }D_0,
 \qquad
 f(X)=s_1^{2-2p}X^{2p},\quad g(X)=1\quad\text{on }D_1.
\]
We show that every polynomial $H$ of actual degree $3\le d\le p$
matches any $f+\lambda g$ on at most $2p+2d$ coordinates.

On a component $uB_0$, set $X=ut$ and divide the matching equation by
its nonzero coefficient of $t^{2p}$. It becomes $t^{2p}=P(t)$, where
$\deg P=d$. If the coefficients of $P$ are not all in $B_0$, a
$B_0$-linear projection $E\to B_0$ annihilating $B_0$ gives a nonzero
polynomial of degree at most $d$ that vanishes at every match. Such a
component contributes at most $d$ matches.

There cannot be coefficientwise $B_0$-valued components in both blocks.
Indeed, if $h_d$ is the leading coefficient of $H$, the normalized leading
coefficients on $v^\epsilon B_0$ and $s_1v^\delta B_0$ have ratio
\[
 s_1^{d-2}v^{(\delta-\epsilon)(d-2p)},
 \qquad \epsilon,\delta\in\{0,1\}.
\]
Modulo $B_0^*$ this equals $s_1^{d-2}v^{(\delta-\epsilon)d}$.
The quotient $E^*/B_0^*$ has order $Q=p^2+1$, the class of $s_1$
generates it, and the class of $v$ has order two. For even $d$ the
exponent is therefore $d-2$ modulo $Q$; for odd $d$ it is $d-2$ or
$d-2+Q/2$. None vanishes, since $1\le d-2\le p-2<Q/2$.

On the at most one block with such a component, the original matching
equation has degree exactly $2p$ in $X$, so that entire block contributes
at most $2p$ matches. The other block contributes at most $2d$ by the two
projection bounds. If neither block has such a component, the bound is
$4d\le2p+2d$. Puncturing cannot increase these counts.

Put $D=\lfloor p/20\rfloor$. Then $2p+2D\le21p/10<T$, since
\[
 5p^2-2-(21p/10+1)^2
   =\frac{59}{100}p^2-\frac{21}{5}p-3>0
\]
for $p\ge257$. All new witnesses thus lie below threshold, preserving
the complete list profile. The old simultaneous quadratic witness gives
the lower source and common-agreement bound $2p$, and the estimate just
proved gives the upper bound. The ratio follows by dividing by $T-k$.
Finally $k\ge4$ and
$na_1(k/n)>\sqrt{kn/2}\ge\sqrt{2n}>T$, proving the placement claim.
\end{proof}
